\documentclass[11pt]{article}
\usepackage{fontspec,amsmath,amssymb,amsthm,hyperref}
\usepackage[margin=1in]{geometry}
\setlength{\emergencystretch}{3em}
\title{Compatibility exports for the shared Hilbert-space coupling proof}
\author{LangevinLean mathematical exposition}
\date{}
\begin{document}\maketitle
\noindent Source module: \texttt{IPLA/Dalalyan.lean}. Compile this proof document with XeLaTeX or LuaLaTeX. The accompanying HTML pairs the mathematical argument with the exact Lean statements and records the scope of the formalization.


This module introduces no new mathematical declarations. It exports the shared \texttt{ULA.Hilbert} recursion, integral reference interface and Dalalyan estimates under the historical IPLA namespace. The imported shared module contains their Lean proofs and parallel mathematical exposition.


\section{The shared synchronous coupling calculation}


The common Euler and diffusion noises cancel in the coupled difference. If \(e_k=\|X_k-Y_k\|_{L^2}\), contraction of the deterministic gradient step and the integrated drift-defect estimate give

\[e_{k+1}\le q e_k+b,\qquad b=\frac{33}{20}Lh\sqrt{hd}.\]

Iteration gives \(e_k\le q^ke_0+b/(1-q)\). In the small-step regime \(0<h\le2/(m+L)\), \(q=1-mh\), yielding bias \((33/20)(L/m)\sqrt{hd}\). The reference structure explicitly requires continuous joined segments, their integral equation, stationary knot laws and score/noise moment bounds. This compatibility file proves none of those requirements anew; the final particle theorem obtains them from the shared Gaussian ULA construction.



The \texttt{export} commands expose the already proved declarations under the IPLA namespace. Their definitions and proofs are unchanged, so the mathematical argument is the one in the imported source modules; the dependency diagram below opens those full proofs. This file itself makes no additional analytic assumption and contains no replacement proof.


\end{document}
